\documentclass[11pt]{article}
\usepackage[margin=1in]{geometry}
\usepackage{amsmath,amssymb,amsfonts}
\usepackage{booktabs}
\usepackage{graphicx}
\usepackage{hyperref}
\usepackage{xcolor}
\usepackage{enumitem}
\usepackage{microtype}
\usepackage{natbib}
\usepackage{multirow}

\bibliographystyle{plainnat}

\title{Two Negative Results for Vector Symbolic Architectures:\\
FFN Replacement and Compositional Image Generation}

\author{
Peter C.\\
Independent Researcher\\
\texttt{peterc3@live.com}
}

\date{}

\begin{document}
\maketitle

\begin{abstract}
Vector Symbolic Architectures (VSAs) offer algebraically structured representations---binding, bundling, permutation---that are $O(D)$ and matmul-free.
These properties make VSA a natural candidate for (1)~replacing dense feed-forward networks with structured retrieval, and (2)~compositional image generation via token binding.
We conduct systematic experiments testing both hypotheses and find that neither works, for different but complementary reasons.

\textbf{Case Study~1: FFN Replacement.}
We test VSA memory layers as replacements for feed-forward networks in Qwen3.6-27B.
The failure mode is a rank bottleneck: VSA's cleanup$\to$bind$\to$retrieve pipeline has effective rank bounded by top-$k$ (typically~4), while FFN mappings are ${\sim}89\%$ linear with effective rank ${>}2048$.
A 164K-parameter rank-16 linear projection captures more variance than a 35M-parameter VSA memory layer.

\textbf{Case Study~2: Compositional Image Generation.}
We test VSA binding for encoding image token sequences from TiTok VQ-8K (64 tokens $\times$ 8192 codebook).
The failure has two independent causes: (a)~FHRR superposition of 64 bindings cannot support 8192-way retrieval ($0\%$ accuracy at all $D$ tested), and (b)~real images do not share token multisets (${<}9\%$ overlap), so the permutation-based generation framing is inapplicable.

\textbf{Positive findings.}
We identify conditions where VSA succeeds: permutation locality in TiTok, 100\% exact permutation recovery at $D{\ge}2048$, and SVD-aligned codebooks capturing 50\% more variance than random.
These successes illuminate the boundary: VSA excels at structured, low-cardinality operations but fails when the target requires high-rank continuous mappings or high-cardinality discrete retrieval.
\end{abstract}

%----------------------------------------------------------------------
\section{Introduction}
\label{sec:intro}

Vector Symbolic Architectures (VSAs), also known as hyperdimensional computing, have attracted renewed interest as a biologically plausible and algebraically structured alternative to standard neural network operations~\citep{kanerva2009hyperdimensional, plate2003holographic}.
VSA represents data as high-dimensional vectors and computes with three core operations: \emph{binding} (element-wise multiplication), \emph{bundling} (superposition via addition), and \emph{permutation} (cyclic shift).
These operations are $O(D)$, matmul-free, and support exact inversion---properties that are absent in standard dense layers.

Two applications stand out as natural fits for VSA's strengths:

\begin{enumerate}[leftmargin=*]
\item \textbf{FFN replacement.} Feed-forward networks consume two-thirds of transformer parameters and are known to store factual knowledge as key-value associations~\citep{sukhbaatar2019augmenting}.
VSA's bind-and-retrieve mechanism is a direct analog of associative memory, suggesting that VSA memory layers could replace dense FFNs with structured retrieval at lower cost.
\item \textbf{Compositional image generation.} Image tokenizers like TiTok~\citep{yu2024image} encode images as sequences of discrete tokens.
VSA binding naturally represents position$\to$content mappings, and permutation operations could support compositional manipulation of token arrangements.
\end{enumerate}

We test both hypotheses with systematic experiments and find that \emph{neither works}, for different but complementary reasons.

\paragraph{FFN replacement} fails due to a \textbf{rank bottleneck}: VSA's top-$k$ retrieval constrains output rank to ${\le}k$, while FFN effective rank exceeds 2048.

\paragraph{Image generation} fails due to a \textbf{capacity bottleneck} (64 superposed bindings cannot support 8192-way retrieval) \emph{and} a \textbf{framing mismatch} (images do not share token multisets, invalidating the permutation premise).

\paragraph{Positive findings} narrow the scope: VSA achieves 100\% exact match for permutation encoding at $D{\ge}2048$ (64-way classification), confirming that VSA works for structured, low-cardinality operations.

\paragraph{Contributions.}
\begin{enumerate}[leftmargin=*]
\item Two independent case studies demonstrating VSA limitations in complementary domains.
\item A rank-bottleneck diagnosis for FFN replacement (theoretical proof sketch + empirical confirmation).
\item Capacity analysis for VSA token binding: theory predicts ${\sim}\log_2(D/N)$ bits per retrieval; experiments confirm.
\item Positive result: VSA permutation encoding achieves 100\% exact match at $D{=}2048$ for 64-position permutations.
\item Negative result: TiTok VQ-8K token multisets have ${<}9\%$ overlap between images, invalidating permutation-based generation.
\end{enumerate}

%----------------------------------------------------------------------
\section{Background}
\label{sec:background}

\subsection{Vector Symbolic Architectures}

We use FHRR (Frequency Holographic Reduced Representation)~\citep{plate2003holographic}, where each symbol is a $D$-dimensional complex phasor $\mathbf{x} = e^{i\boldsymbol{\theta}}$ with $\boldsymbol{\theta} \in [-\pi, \pi)^D$.
Binding is element-wise complex multiplication: $\text{bind}(\mathbf{a}, \mathbf{b}) = \mathbf{a} \odot \mathbf{b}$.
Bundling is complex addition followed by normalization: $\text{bundle}(\mathbf{a}, \mathbf{b}) = \text{normalize}(\mathbf{a} + \mathbf{b})$.
Unbinding uses the conjugate: $\text{unbind}(\mathbf{c}, \mathbf{a}) = \mathbf{c} \odot \mathbf{a}^*$.

The storage capacity for $N$ superposed bindings is approximately $\log_2(D/N)$ bits per retrieval~\citep{frady2018theory}.
This means that with $N{=}64$ bindings and $D{=}4096$, each retrieval can distinguish among at most $2^6 = 64$ items---far below the 8192-entry codebook ($2^{13}$) needed for image token retrieval.

\subsection{Feed-Forward Networks in Transformers}

Standard transformer FFNs consist of an up-projection, nonlinear activation, and down-projection: $\text{FFN}(\mathbf{x}) = W_2 \cdot \sigma(W_1 \mathbf{x})$.
These layers contain $2 \times d_\text{model} \times d_\text{ffn}$ parameters per layer and are known to store factual knowledge~\citep{sukhbaatar2019augmenting}.

\subsection{Image Tokenization with TiTok}

TiTok~\citep{yu2024image} is a 1D image tokenizer that encodes a $256{\times}256$ image into 64 discrete tokens from an 8192-entry codebook ($64$ dimensions per entry).
We use the frozen pre-trained \texttt{titok\_bl64\_vq8k\_imagenet} model (Apache~2.0 license) throughout.

%----------------------------------------------------------------------
\section{Case Study 1: VSA for FFN Replacement}
\label{sec:ffn}

\subsection{Setup}

We extract activation pairs $(x_\text{in}, x_\text{out})$ from Qwen3.6-27B~\citep{qwen2024} ($d_\text{model}{=}5120$, 64 layers) at layers~3, 27, and~43.
We sample 50{,}000 tokens from FineWeb-Edu with 5\% held out for validation.
The metric is \emph{variance captured}: $\text{Var}\% = (1 - \text{MSE}_\text{val} / \text{MSE}_\text{zero}) \times 100\%$, where $\text{MSE}_\text{zero}$ is the mean squared error of predicting zero.

\subsection{Architectures Tested}

\begin{itemize}[leftmargin=*]
\item \textbf{VSA V1} (frozen random codebooks): standard FHRR bind-retrieve with random position and content phasors.
\item \textbf{VSA V2} (learned + Gumbel-softmax): codebook entries and routing learned end-to-end.
\item \textbf{VSA-MoE}: VSA routing selects experts in a mixture-of-experts architecture.
\item \textbf{Baselines}: rank-$r$ linear (SVD), rank-$r$ + MLP, learned-gate MoE.
\end{itemize}

\subsection{Results}

Table~\ref{tab:arch} shows the architecture comparison at layer~27.
The key finding: a 164K-parameter rank-16 linear projection (5.9\%) \emph{outperforms} a 35M-parameter VSA memory layer (4.8\%).
Rank-2048 linear (36.6\%) dominates all VSA variants by $7\times$ or more.

\begin{table}[h]
\centering
\caption{Architecture comparison (layer~27, Qwen3.6-27B).}
\label{tab:arch}
\begin{tabular}{lcc}
\toprule
Architecture & Var\% & Params \\
\midrule
Zero baseline & 0.0 & 0 \\
Cube Memory V1 (frozen VSA) & ${\sim}5$ & 35M \\
Cube Memory V2 (learned VSA) & 4.8 & 35M \\
Rank-16 linear (SVD) & 5.9 & 164K \\
VSA-MoE 16$\times$128 top-4 & 14.2 & 24M \\
Learned-MoE 8$\times$256 top-2 & 16.2 & 21M \\
Rank-2048 linear & 36.6 & 21M \\
\textbf{Rank-2048 + MLP-512} & \textbf{38.4} & \textbf{26M} \\
Full-rank ceiling & 41.1 & 52M \\
\bottomrule
\end{tabular}
\end{table}

\subsection{Ablation Studies}

\paragraph{Top-$k$ scaling.}
Increasing top-$k$ from 4 to 64 does not break the ceiling: variance captured decreases slightly from 2.9\% to 2.7\%.
This rules out the hypothesis that the bottleneck is insufficient retrieval breadth.

\paragraph{Codebook ablation.}
SVD-aligned codebooks improve over random (4.3\% vs 2.8\%), confirming that aligning codebook entries with the data manifold helps.
However, the improvement does not change the conclusion: all VSA variants remain far below linear baselines.

\paragraph{FLOPs comparison.}
All FLOPs are inference-only forward pass.
Rank-2048 linear achieves 36.6\% variance at 42M FLOPs/token.
VSA V2 achieves 4.8\% at 62M FLOPs/token---\emph{worse and slower}.

\paragraph{Rank-equalized comparison.}
At matched effective rank~4, a rank-4 linear (SVD) yields negative $R^2$ (worse than predicting the mean) at 41K FLOPs, while VSA top-4 achieves 2.9\% at 62M FLOPs.
VSA's marginal advantage comes from data-dependent routing (selecting different codebook entries per input), but both are catastrophically below rank-2048 linear (36.6\%).
VSA spends $1514\times$ more FLOPs for this marginal advantage.

\subsection{Diagnosis: The Rank Bottleneck}

The VSA retrieval output is:
\begin{equation}
\mathbf{y} = \sum_{j \in \text{top-}k} \alpha_j \cdot \mathbf{v}_j,
\quad \alpha_j = \text{softmax}(\text{sim}_j),
\end{equation}
where $\mathbf{v}_j$ are memory value vectors. This is a convex combination of $k$~vectors, so $\mathbf{y} \in \text{span}\{\mathbf{v}_1, \ldots, \mathbf{v}_k\}$, giving $\text{rank}(\text{output}) \le k$.

With top-$k{=}4$, VSA output rank is at most~4.
The FFN effective rank exceeds~2048 (singular values decay slowly; see Appendix).
Retrieval is fundamentally underpowered: the FFN mapping is ${\sim}89\%$ linear, and the remaining ${\sim}11\%$ is better captured by a small MLP than by sparse retrieval.

This rank bottleneck applies to \emph{all} retrieval-based FFN replacements---Product Key Memory~\citep{lample2019large}, holographic reduced representations, or any top-$k$ selection mechanism.

%----------------------------------------------------------------------
\section{Case Study 2: VSA for Image Generation}
\label{sec:imgen}

\subsection{Motivation: Permutation-Group Indexed Generation}

We hypothesize that image generation can be framed as token permutation manipulation.
TiTok encodes images as 64 discrete tokens from an 8192-entry codebook.
VSA binding naturally represents position$\to$content mappings: $\mathbf{h} = \sum_{i=1}^{64} \text{bind}(\mathbf{p}_i, \mathbf{c}_{t_i})$, where $\mathbf{p}_i$ is a position phasor and $\mathbf{c}_{t_i}$ is a content phasor for token $t_i$.

\subsection{Prerequisite: Permutation Locality (Experiment~0)}

We first validate that the TiTok decoder is smooth with respect to token permutation.
We encode a test image, apply $N$ random position swaps, decode, and measure MSE against the original.

\begin{table}[h]
\centering
\caption{Permutation locality: MSE vs.\ number of random token swaps.}
\label{tab:permlocality}
\begin{tabular}{cccccc}
\toprule
Swaps & 0 & 1 & 4 & 16 & 64 \\
\midrule
MSE & 0.000 & 0.001 & 0.022 & 0.073 & 0.153 \\
\bottomrule
\end{tabular}
\end{table}

The monotonic relationship (Table~\ref{tab:permlocality}) confirms the generation premise has empirical support: small permutations yield small visual changes.

\subsection{Token Binding Fails (Experiments~1, 1b, 1c, 2)}

\paragraph{Experiment~1: Pure VSA binding.}
We bind 64 position phasors with content phasors, superpose, unbind each position, and classify against the 8192-entry codebook.
Theory predicts ${\sim}\log_2(D/64)$ bits per retrieval; 8192 entries need 13~bits.
Even $D{=}4096$ provides only 6~bits.
Result: 0\% accuracy at all $D$ tested (Table~\ref{tab:tokenbinding}).

\begin{table}[h]
\centering
\caption{Pure VSA token binding: accuracy vs.\ dimension.}
\label{tab:tokenbinding}
\begin{tabular}{lcl}
\toprule
$D$ & Acc\% & Bits (have / need) \\
\midrule
512 & 0.0 & 3.0 / 13.0 \\
1024 & 0.0 & 4.0 / 13.0 \\
2048 & 0.0 & 5.0 / 13.0 \\
4096 & 0.0 & 6.0 / 13.0 \\
\bottomrule
\end{tabular}
\end{table}

\paragraph{Experiment~1b: VSA + MLP decoder.}
Adding a 2-layer MLP after unbinding produces catastrophic overfitting at all $D$: training loss approaches zero while validation loss reaches~25.
The decoder memorizes training-specific noise patterns.
An MLP-only control (position one-hot $\to$ MLP $\to$ 8192-way, no VSA) achieves 0.93\% accuracy, confirming the task is impossible from positional information alone.

\paragraph{Experiment~1c: Factored codebook.}
We decompose $8192 = 128 \times 64$ and use 3-way binding ($\text{pos} \otimes \text{factA} \otimes \text{factB}$), retrieving factors independently.
A hard ceiling appears at ${\sim}5.8\%$ for the 64-way factor regardless of $D$ (Table~\ref{tab:factored}).
The bottleneck is structural (cross-talk from superposition), not capacity.

\begin{table}[h]
\centering
\caption{Factored codebook ($128\times64$): accuracy per factor.}
\label{tab:factored}
\begin{tabular}{lccc}
\toprule
$D$ & AccA\% (128-way) & AccB\% (64-way) & Joint\% \\
\midrule
512 & 2.2 & 5.6 & 0.1 \\
1024 & 2.3 & 5.8 & 0.2 \\
2048 & 2.0 & 5.7 & 0.1 \\
4096 & 1.9 & 5.9 & 0.1 \\
\bottomrule
\end{tabular}
\end{table}

\paragraph{Experiment~2: Real TiTok tokens.}
We test on 6000 Imagenette images encoded by frozen TiTok.
Token distribution entropy is 12.85 / 13.00~bits; all 8192 codes are used.
At $D{=}4096$, accuracy plateaus at 0.61\% after 1500 steps---the same plateau as $D{=}512$.
Near-uniform codebook usage means real tokens are as hard as random.

\subsection{Permutation Encoding Works (Experiment~3)}

We reframe: encode \emph{permutations} (64-way classification) instead of \emph{tokens} (8192-way).
This is a pure VSA theory test requiring no training.
For each permutation $\sigma$, we encode $\mathbf{h}_\sigma = \sum_{i=1}^{64} \text{bind}(\mathbf{p}_i, \mathbf{p}_{\sigma(i)})$ and decode by unbinding each source position and classifying among 64 destinations.

\begin{table}[h]
\centering
\caption{Permutation recovery: exact match \% ($n{=}2000$ per condition).}
\label{tab:perm}
\begin{tabular}{lcccccc}
\toprule
$D$ & $k{=}1$ & $k{=}2$ & $k{=}4$ & $k{=}8$ & $k{=}16$ & $k{=}32$ \\
\midrule
256 & 0.0 & 0.0 & 0.0 & 0.0 & 0.0 & 0.3 \\
512 & 0.0 & 0.0 & 0.0 & 0.0 & 0.0 & 67.2 \\
1024 & 89.3 & 76.3 & 54.6 & 29.6 & 15.3 & 99.9 \\
2048 & 99.8 & 99.8 & 99.8 & 98.6 & 96.5 & \textbf{100.0} \\
4096 & \textbf{100.0} & \textbf{100.0} & \textbf{100.0} & \textbf{100.0} & \textbf{100.0} & \textbf{100.0} \\
\bottomrule
\end{tabular}
\end{table}

At $D{\ge}2048$, recovery is perfect (0 failures in 2000 test permutations per condition, 95\% CI: [99.8\%, 100\%]).
The 64-way classification is well within VSA capacity, confirming that VSA excels at structured operations on small discrete sets.

\subsection{But Images Aren't Permutations (Experiment~4)}

We test whether real images can be approximated as permutations of a cluster reference.
For each image: find nearest cluster (by token histogram similarity, $k$-means), solve optimal assignment via the Hungarian algorithm~\citep{kuhn1955hungarian}, and measure token match rate.

\begin{table}[h]
\centering
\caption{Permutation reconstruction: can images be expressed as permutations of references?}
\label{tab:recon}
\begin{tabular}{lccc}
\toprule
$K$ clusters & Within-cluster overlap & Match rate & ${\ge}50\%$ match \\
\midrule
10 & 4.0\% & 2.7\% & 0.2\% \\
50 & 2.1\% & 5.2\% & 0.8\% \\
200 & 2.4\% & 8.9\% & 3.3\% \\
\bottomrule
\end{tabular}
\end{table}

``Within-cluster overlap'' is the average pairwise multiset intersection size divided by 64 between random pairs in the same cluster.
``Match rate'' is the fraction of positions where the Hungarian-optimal permutation of the reference yields the correct token.

Each image uses $64/8192 = 0.78\%$ of the codebook.
Under uniform usage, the expected multiset overlap between two random images is ${\sim}0.5$ tokens (birthday collision rate: $8192 \times (64/8192)^2 \approx 0.5$).
The permutation framing requires near-complete multiset overlap, which is impossible when the codebook is large and near-uniformly used.

\subsection{Diagnosis: Two Independent Failure Modes}

\begin{enumerate}[leftmargin=*]
\item \textbf{Capacity failure.} FHRR superposition of $N$ bindings provides ${\sim}\log_2(D/N)$ bits per retrieval. With $N{=}64$ and codebook size 8192 (13~bits), no feasible $D$ suffices. Reducing cardinality to 64-way (permutations) solves this.
\item \textbf{Framing failure.} Even if token retrieval worked, permutation-based generation requires shared token multisets between images. With 8192 near-uniformly-used codebook entries, this assumption fails catastrophically (${<}9\%$ overlap).
\end{enumerate}

%----------------------------------------------------------------------
\section{Positive Findings and Boundary Conditions}
\label{sec:positive}

\subsection{Where VSA Succeeds}

\begin{enumerate}[leftmargin=*]
\item \textbf{Permutation representation.} 100\% exact recovery at $D{=}2048$ for 64-position permutations (Experiment~3). VSA is ideal for structured operations on small discrete sets.
\item \textbf{Permutation locality.} The TiTok decoder is smooth with respect to token permutation (Experiment~0). This property is real and potentially useful for local image editing.
\item \textbf{SVD-aligned codebooks.} Initializing VSA codebooks from SVD of the target activation matrix captures 50\% more variance than random (4.3\% vs 2.8\%).
\end{enumerate}

\subsection{The Boundary}

VSA succeeds when classification cardinality is low (64 positions, not 8192 codes), the target has discrete compositional structure, and operations are group-theoretic.
VSA fails when the target mapping is high-rank and continuous, retrieval requires high-cardinality discrimination from superposition, or the compositional framing does not match the data.

\subsection{Generalization}

The rank bottleneck (Case Study~1) applies to \emph{all} retrieval-based FFN replacements: Product Key Memory~\citep{lample2019large}, Memory Layers~\citep{wu2024memory}, or any top-$k$ selection mechanism.
The output of any top-$k$ retrieval is constrained to a $k$-dimensional subspace, which is insufficient for the high-rank mappings that FFNs implement.

The capacity limit (Case Study~2) applies to any VSA scheme that superposes $N$ bindings and retrieves from codebook size $C$: reliable retrieval requires $D \gg N \times C$.

%----------------------------------------------------------------------
\section{Related Work}
\label{sec:related}

\paragraph{VSA and Hyperdimensional Computing.}
The theoretical foundations of VSA trace to Kanerva's sparse distributed memory~\citep{kanerva2009hyperdimensional} and Plate's holographic reduced representations~\citep{plate2003holographic}.
\citet{frady2018theory} established capacity bounds for FHRR, which our experiments confirm empirically.

\paragraph{Memory Layers.}
Product Key Memory~\citep{lample2019large} and recent scaling work~\citep{wu2024memory, sukhbaatar2019augmenting} use key-value retrieval to augment or replace FFN layers.
Our rank bottleneck analysis applies to these methods whenever they use top-$k$ selection.

\paragraph{FFN Compression.}
Low-rank factorization~\citep{hsu2022language}, pruning~\citep{frantar2023sparsegpt}, and distillation~\citep{hinton2015distilling} are alternatives to memory-based replacement.
Our results show that even simple rank-2048 linear projections dramatically outperform structured retrieval approaches.

\paragraph{Image Tokenization.}
VQ-VAE~\citep{van2017neural}, VQGAN~\citep{esser2021taming}, and TiTok~\citep{yu2024image} have established discrete tokenization as a viable representation for image generation.
Our permutation locality finding (Experiment~0) adds a new characterization of TiTok's latent space.

\paragraph{Sparse MoE.}
Mixture-of-experts approaches~\citep{shazeer2017outrageously, fedus2022switch, jiang2024mixtral} use routing to select among specialized sub-networks.
Our VSA-MoE variant (14.2\%) underperforms learned-gate MoE (16.2\%), suggesting that VSA routing provides no advantage over standard gating.

%----------------------------------------------------------------------
\section{Conclusion}
\label{sec:conclusion}

We tested VSA in two settings where its algebraic properties should provide advantages: structured FFN replacement and compositional image generation.
Both fail, for different reasons:

\begin{enumerate}[leftmargin=*]
\item \textbf{FFN replacement: the rank bottleneck.} FFN mappings are ${\sim}89\%$ linear with effective rank ${>}2048$. VSA retrieval collapses rank to top-$k$. Linear projections dominate at every parameter budget.
\item \textbf{Image generation: capacity bottleneck plus framing mismatch.} Superposition of 64 bindings provides ${\sim}6$ bits per retrieval; 8192 codebook entries need 13~bits. And even with perfect retrieval, images do not share token multisets, so the permutation framing fails.
\end{enumerate}

The positive findings narrow VSA's useful scope: it excels at representing discrete group operations (permutations of 64 positions: 100\% exact match) but fails when the target requires either high-rank continuous mappings or high-cardinality discrete retrieval.

These results suggest that VSA's future in deep learning lies in genuinely compositional, low-cardinality tasks---symbolic reasoning, relational learning, discrete program synthesis---rather than as a drop-in replacement for dense continuous computations.

All code and data are available at \url{https://github.com/Peterc3-dev/cube-memory}.

%----------------------------------------------------------------------
\bibliography{references}

%----------------------------------------------------------------------
\appendix
\section{Experimental Details}
\label{app:details}

\paragraph{Hardware.}
All Case Study~2 experiments ran on an AMD Ryzen~AI~9 HX~370 APU with 23\,GB RAM (no discrete GPU).
Case Study~1 activation extraction used a ThinkCentre M70q Gen~5 with 32\,GB RAM.
Total compute: approximately 8~hours across both machines.

\paragraph{Hyperparameters.}
VSA experiments used AdamW with learning rate $10^{-3}$, weight decay $0.01$, cosine annealing, batch size 16 (64 for FFN distillation), and 2000--3000 training steps.
Position phasors were frozen random; content phasors were learned.

\paragraph{Reproducibility.}
All experiments use fixed random seeds (42).
The TiTok model is the publicly available \texttt{yucornetto/tokenizer\_titok\_bl64\_vq8k\_imagenet} checkpoint.
Imagenette~320 was used for image tokenization (6000 images: 5000 train, 1000 val).

\end{document}
